\chapter{Úvod}
\label{chap:introduction}

V~posledních letech lze na evropském potravinářském trhu pozorovat výrazný posun spotřebitelských preferencí směrem ke zdravějším alternativám tradičních sladkostí. Tento trend, poháněný rostoucím povědomím o~výživě a~důrazem na kvalitu surovin, se projevuje mimo jiné v~rostoucí oblibě produktů kombinujících ovoce a~čokoládu (Anselmsson et al., 2014). Jedním z~produktů, který na tomto trendu úspěšně staví, je argentinská značka Franui, nabízející mražené maliny obalené ve dvou vrstvách čokolády.

Značka Franui vstoupila na evropský trh ve velkém měřítku v~roce 2021 a~rychle si vybudovala silnou pozici, a~to zejména díky virálnímu marketingu na sociálních sítích TikTok a~Instagram (Franui, 2024). Na českém trhu se Franui stalo nejvýraznějším zástupcem kategorie mraženého ovoce v~čokoládě, přičemž přímá lokální konkurence v~současnosti prakticky neexistuje.

\section{Identifikace problému}

Při analýze dostupné nabídky na českém trhu byla identifikována výrazná mezera — absence lokálního produktu, který by mohl konkurovat zahraničnímu Franui. Zatímco český spotřebitel má přístup k~široké škále tradičních mražených pochutin (jako je zmrzlina Míša či nanuk Ledňáček), segment mraženého ovoce v~prémiové čokoládě zůstává obsazen výhradně importovaným produktem s~relativně vysokou cenou (přibližně 160~Kč za 150g balení).

Tato situace představuje potenciální příležitost pro uvedení nového českého produktu, který by mohl nabídnout konkurenceschopnou alternativu s~důrazem na lokální původ surovin, \gls{bio} kvalitu a~dostupnější cenovou hladinu.

\section{Cíle práce}

Hlavním cílem této diplomové práce je provést komplexní marketingový výzkum, který posoudí potenciál uvedení nové značky mraženého ovoce v~čokoládě \textbf{Berrie} na český trh a~identifikuje klíčové faktory ovlivňující rozhodování spotřebitelů o~koupi v~této produktové kategorii.

Dílčí cíle práce jsou:
\begin{enumerate}
  \item zjistit, jak čeští spotřebitelé vnímají kategorii mraženého ovoce v~čokoládě a~jaké asociace si s~ní spojují,
  \item porozumět motivacím a~bariérám ke koupi tohoto typu produktu,
  \item ověřit chuťové preference a~očekávání týkající se složení produktu prostřednictvím senzorické analýzy,
  \item zjistit cenovou citlivost spotřebitelů a~stanovit optimální cenovou hladinu,
  \item posoudit, jakou roli hraje lokální původ surovin a~vnímání zdravějších alternativ sladkostí při nákupním rozhodování.
\end{enumerate}

\section{Výzkumné otázky}

Na základě stanovených cílů jsou formulovány následující výzkumné otázky:
\begin{enumerate}
  \item Jaké faktory (chuť, textura, složení, značka, původ surovin) mají největší vliv na rozhodnutí spotřebitele produkt vyzkoušet?
  \item Jaké jsou hlavní rozdíly ve vnímání produktu mezi mladými spotřebiteli (generace~Z, mileniálové) a~rodiči s~dětmi?
  \item Jaký je cenový strop produktu v~cílových skupinách a~jak se vztahuje k~ceně konkurenční značky Franui?
  \item Jaký je potenciál lokální alternativy vůči zavedenému importovanému produktu?
\end{enumerate}

\section{Struktura práce}

Práce je rozdělena do šesti hlavních kapitol. Po úvodu následuje teoretická část (Kapitola~\ref{chap:theory}), která shrnuje relevantní poznatky z~oblasti marketingového výzkumu, spotřebitelského chování, senzorické analýzy potravin a~cenové citlivosti. Kapitola~\ref{chap:methodology} popisuje metodologii výzkumu, včetně designu kvalitativní a~kvantitativní fáze. Kapitola~\ref{chap:research} představuje vlastní výzkumnou část, zahrnující definici produktu, analýzu konkurence a~realizaci výzkumu. Výsledky a~jejich diskuse jsou předmětem Kapitoly~\ref{chap:results}. Závěrečná Kapitola~\ref{chap:conclusion} shrnuje hlavní zjištění a~formuluje doporučení pro marketingovou strategii značky Berrie.
